\documentclass[11pt]{article}

\usepackage[a4paper,margin=1in]{geometry}
\usepackage{amsmath,amssymb,amsthm,mathtools}
\usepackage[T1]{fontenc}
\usepackage{lmodern}
\usepackage{microtype}
\usepackage{hyperref}

\newtheorem{theorem}{Theorem}[section]
\newtheorem{lemma}[theorem]{Lemma}
\newtheorem{proposition}[theorem]{Proposition}
\newtheorem{corollary}[theorem]{Corollary}
\newtheorem{remark}[theorem]{Remark}
\newtheorem{definition}[theorem]{Definition}

\DeclareMathOperator{\ind}{ind}
\DeclareMathOperator{\diag}{diag}
\DeclareMathOperator{\tridiag}{tridiag}
\DeclareMathOperator{\rank}{rank}

\newcommand{\R}{\mathbb{R}}
\newcommand{\C}{\mathbb{C}}

\title{Simple Connectedness for the Clean Hatano--Nelson Family\\
with Open Boundary Conditions}
\author{}
\date{}

\begin{document}
\maketitle

\section{Setup}

Fix \(n\ge 2\), \(0<a<1\), and consider the real nonnormal tridiagonal Toeplitz matrix
\[
\mathbf A_n(a):=\tridiag(a,0,1)
=
\begin{bmatrix}
0 & 1 \\
a & 0 & \ddots \\
& \ddots & \ddots & \ddots \\
&& \ddots & 0 & 1 \\
&&& a & 0
\end{bmatrix}\in \R^{n\times n}.
\]

For \(\varepsilon>0\), its \(\varepsilon\)-pseudospectrum is
\[
\sigma_\varepsilon(\mathbf A_n(a))
:=
\left\{
z\in\C:\;
\|(z\mathbf I_n-\mathbf A_n(a))^{-1}\|_2>\varepsilon^{-1}
\right\},
\]
with the usual convention that the resolvent norm is \(+\infty\) on the spectrum.
Equivalently,
\[
\sigma_\varepsilon(\mathbf A_n(a))
=
\left\{
z\in\C:\;
s_{\min}(z\mathbf I_n-\mathbf A_n(a))<\varepsilon
\right\}.
\]

\begin{lemma}[Real spectrum]\label{lem:real-spectrum}
Let
\[
\mathbf D_n:=\diag(1,q,q^2,\dots,q^{n-1}),
\qquad
q:=\sqrt a.
\]
Then
\[
\mathbf D_n^{-1}\mathbf A_n(a)\mathbf D_n
=
\sqrt a\,\tridiag(1,0,1).
\]
Consequently,
\[
\sigma(\mathbf A_n(a))
=
\left\{
2\sqrt a \cos\frac{j\pi}{n+1}:j=1,\dots,n
\right\}
\subset \R.
\]
\end{lemma}

\begin{proof}
A direct computation gives
\[
(\mathbf D_n^{-1}\mathbf A_n(a)\mathbf D_n)_{j,j+1}=q=\sqrt a,
\qquad
(\mathbf D_n^{-1}\mathbf A_n(a)\mathbf D_n)_{j+1,j}=aq^{-1}=\sqrt a.
\]
Thus \(\mathbf D_n^{-1}\mathbf A_n(a)\mathbf D_n=\sqrt a\,\tridiag(1,0,1)\), and the spectral formula is the classical one for the symmetric tridiagonal Toeplitz matrix.
\end{proof}

\begin{remark}[Two endpoint models]\label{rem:endpoints}
The continuous endpoint families \(a=1\) and \(a=0\) are especially transparent.

\begin{enumerate}
\item If \(a=1\), then \(\mathbf A_n(1)=\tridiag(1,0,1)\) is symmetric, hence normal, and
\[
\sigma_\varepsilon(\mathbf A_n(1))
=
\bigcup_{j=1}^n
D\!\left(2\cos\frac{j\pi}{n+1},\varepsilon\right).
\]
Therefore every connected component is simply connected.

\item If \(a=0\), then \(\mathbf A_n(0)\) is the nilpotent Jordan block with superdiagonal \(1\). For every \(\theta\in\R\),
\[
\mathbf U_\theta^*\,\mathbf A_n(0)\,\mathbf U_\theta
=
e^{i\theta}\mathbf A_n(0),
\qquad
\mathbf U_\theta:=\diag(1,e^{i\theta},\dots,e^{i(n-1)\theta}),
\]
so \(\sigma_\varepsilon(\mathbf A_n(0))\) is rotationally invariant. Since \(0\) is the only eigenvalue, the pseudospectrum has only one connected component; rotational invariance then forces it to be a disc centered at \(0\), hence simply connected.
\end{enumerate}

These two limit cases motivate the low-\(a\) and near-normal regimes used in the final topological argument.
\end{remark}

\begin{remark}[Reduction of the general family]\label{rem:general-b}
If one starts with
\[
\mathbf A_n=\tridiag(\alpha,0,\beta),
\qquad
0<\alpha<\beta,
\]
then
\[
\mathbf A_n=\beta\,\mathbf A_n(\eta),
\qquad
\eta:=\frac{\alpha}{\beta}\in(0,1),
\]
and
\[
\sigma_\varepsilon(\mathbf A_n)
=
\beta\,\sigma_{\varepsilon/\beta}(\mathbf A_n(\eta)).
\]
Thus simple connectedness for the normalized family \(\mathbf A_n(a)=\tridiag(a,0,1)\) is equivalent to simple connectedness for the general clean Hatano--Nelson family with \(0<\alpha<\beta\).
\end{remark}

\section{An algebraic boundary polynomial}

For \(z=x+iy\), define the Hermitian matrix
\[
\mathbf H_{n,\varepsilon}(x,y;a)
:=
(z\mathbf I_n-\mathbf A_n(a))^*(z\mathbf I_n-\mathbf A_n(a))
-\varepsilon^2\mathbf I_n,
\]
and its determinant
\[
F_{n,\varepsilon}(x,y;a)
:=
\det \mathbf H_{n,\varepsilon}(x,y;a).
\]

\begin{proposition}[Boundary polynomial]\label{prop:boundary-polynomial}
For every \(n\ge 2\), \(0<a<1\), and \(\varepsilon>0\), the polynomial \(F_{n,\varepsilon}(x,y;a)\) belongs to \(\R[a,x^2,y^2]\) and has total degree \(2n\) in \((x,y)\). Moreover,
\[
\partial \sigma_\varepsilon(\mathbf A_n(a))
=
\left\{
x+iy\in\C:\;
F_{n,\varepsilon}(x,y;a)=0,\;
\mathbf H_{n,\varepsilon}(x,y;a)\succeq 0
\right\}.
\]
Hence
\[
\partial \sigma_\varepsilon(\mathbf A_n(a))
\subseteq
\left\{
x+iy\in\C:\;
F_{n,\varepsilon}(x,y;a)=0
\right\}.
\]
\end{proposition}

\begin{proof}
Since the entries of \(\mathbf H_{n,\varepsilon}(x,y;a)\) are polynomial in \((a,x,y)\), the determinant \(F_{n,\varepsilon}\) is polynomial in these variables. Because \(\mathbf A_n(a)\) is real,
\[
\mathbf H_{n,\varepsilon}(x,-y;a)
=
\overline{\mathbf H_{n,\varepsilon}(x,y;a)},
\]
and therefore \(F_{n,\varepsilon}(x,-y;a)=F_{n,\varepsilon}(x,y;a)\).

Next, let
\[
\mathbf\Sigma_n:=\diag(1,-1,1,-1,\dots,(-1)^{n-1}).
\]
Then \(\mathbf\Sigma_n\mathbf A_n(a)\mathbf\Sigma_n^{-1}=-\mathbf A_n(a)\), so
\[
\mathbf\Sigma_n\mathbf H_{n,\varepsilon}(x,y;a)\mathbf\Sigma_n^{-1}
=
\mathbf H_{n,\varepsilon}(-x,-y;a).
\]
Taking determinants yields
\[
F_{n,\varepsilon}(-x,-y;a)=F_{n,\varepsilon}(x,y;a),
\]
and together with the evenness in \(y\) this implies evenness in \(x\). Hence
\[
F_{n,\varepsilon}(x,y;a)\in \R[a,x^2,y^2].
\]

The degree statement is proved exactly as in the standard determinant expansion: the identity permutation contributes \((x^2+y^2)^n\), while every nonidentity permutation uses at least two off-diagonal entries and therefore lowers the total degree by at least \(2\). Thus \(\deg_{x,y}F_{n,\varepsilon}=2n\).

Now set
\[
f(z):=s_{\min}(z\mathbf I_n-\mathbf A_n(a)).
\]
Since
\[
\sigma_\varepsilon(\mathbf A_n(a))
=
\{z\in\C:\;f(z)<\varepsilon\},
\]
we always have
\[
\partial \sigma_\varepsilon(\mathbf A_n(a))
\subseteq
\{z\in\C:\;f(z)=\varepsilon\}.
\]
Conversely, if \(f(z_0)=\varepsilon\) and \(z_0\notin \partial \sigma_\varepsilon(\mathbf A_n(a))\), then \(z_0\) lies in the interior of the complement, and the resolvent norm
\[
\nu(z):=\|(z\mathbf I_n-\mathbf A_n(a))^{-1}\|
\]
attains a local maximum \(\varepsilon^{-1}\) at \(z_0\). Since \(\nu\) is subharmonic on the resolvent set, it must be locally constant; hence \(f(z)\equiv\varepsilon\) on an open set. Then \(F_{n,\varepsilon}\) vanishes on an open set, contradicting that it is a nonzero polynomial. Therefore
\[
\partial \sigma_\varepsilon(\mathbf A_n(a))
=
\{z\in\C:\;f(z)=\varepsilon\}.
\]

Finally, the eigenvalues of \(\mathbf H_{n,\varepsilon}(x,y;a)\) are \(s_j(z\mathbf I_n-\mathbf A_n(a))^2-\varepsilon^2\). Hence
\[
f(z)=\varepsilon
\iff
\lambda_{\min}\bigl(\mathbf H_{n,\varepsilon}(x,y;a)\bigr)=0
\iff
F_{n,\varepsilon}(x,y;a)=0
\ \text{and}\
\mathbf H_{n,\varepsilon}(x,y;a)\succeq 0.
\]
\end{proof}

\begin{proposition}[No isolated boundary points]\label{prop:no-isolated}
For every \(n\ge 2\), \(0<a<1\), and \(\varepsilon>0\), the boundary \(\partial \sigma_\varepsilon(\mathbf A_n(a))\) has no isolated points.
\end{proposition}

\begin{proof}
This is a purely topological fact for boundaries of nonempty proper open subsets of \(\C\): if \(z_0\in \partial\Omega\) were isolated in \(\partial\Omega\), then a small punctured disc around \(z_0\) would split into two disjoint nonempty relatively open sets \(\Omega\cap(B\setminus\{z_0\})\) and \(\Omega^c\cap(B\setminus\{z_0\})\), contradicting connectedness of the punctured disc.
\end{proof}

\begin{remark}[Discarded branches]\label{rem:discarded-branches}
The full real algebraic curve \(F_{n,\varepsilon}(x,y;a)=0\) is generally larger than the true boundary. Indeed, the equation
\[
F_{n,\varepsilon}(x,y;a)=0
\]
merely says that \emph{some} singular value of \(z\mathbf I_n-\mathbf A_n(a)\) equals \(\varepsilon\). The true boundary is the branch where the \emph{smallest} singular value equals \(\varepsilon\).
\end{remark}

\begin{proposition}[Discarded algebraic branches lie inside]\label{prop:discarded-inside}
If
\[
F_{n,\varepsilon}(x,y;a)=0
\qquad\text{and}\qquad
\mathbf H_{n,\varepsilon}(x,y;a)\not\succeq 0,
\]
then \(x+iy\in \sigma_\varepsilon(\mathbf A_n(a))\).
\end{proposition}

\begin{proof}
The matrix \(\mathbf H_{n,\varepsilon}(x,y;a)\) is Hermitian. If it is singular and not positive semidefinite, then its smallest eigenvalue is negative. Hence
\[
s_{\min}(z\mathbf I_n-\mathbf A_n(a))^2-\varepsilon^2<0,
\]
so \(s_{\min}(z\mathbf I_n-\mathbf A_n(a))<\varepsilon\), i.e. \(z\in \sigma_\varepsilon(\mathbf A_n(a))\).
\end{proof}

\begin{proposition}[Explicit form of the Hermitian pencil]\label{prop:explicit-H}
Set
\[
\tau:=x^2+y^2-\varepsilon^2,
\qquad
c:=az+\overline z=(1+a)x-i(1-a)y,
\qquad
p:=a,
\qquad
d:=\tau+a^2+1.
\]
Then
\[
\mathbf H_{n,\varepsilon}(x,y;a)
=
(|z|^2-\varepsilon^2)\mathbf I_n-\overline z\,\mathbf A_n(a)-z\,\mathbf A_n(a)^{\mathsf T}
+\mathbf A_n(a)^{\mathsf T}\mathbf A_n(a),
\]
and its entries are
\[
\bigl(\mathbf H_{n,\varepsilon}(x,y;a)\bigr)_{jk}
=
\begin{cases}
\tau+a^2, & j=k=1,\\[1mm]
\tau+1, & j=k=n,\\[1mm]
d, & j=k\in\{2,\dots,n-1\},\\[1mm]
-c, & k=j+1,\\[1mm]
-\overline c, & j=k+1,\\[1mm]
a, & |j-k|=2,\\[1mm]
0, & |j-k|>2.
\end{cases}
\]
\end{proposition}

\begin{proof}
The first identity is just the product expansion
\[
(\overline z\mathbf I_n-\mathbf A_n(a)^{\mathsf T})(z\mathbf I_n-\mathbf A_n(a))-\varepsilon^2\mathbf I_n.
\]
The description of the entries follows from the facts that:
\begin{itemize}
\item \(-\overline z\,\mathbf A_n(a)-z\,\mathbf A_n(a)^{\mathsf T}\) contributes the first off-diagonals \(-c\) and \(-\overline c\);
\item \(\mathbf A_n(a)^{\mathsf T}\mathbf A_n(a)\) has diagonal entries \(a^2,a^2+1,\dots,a^2+1,1\);
\item \(\mathbf A_n(a)^{\mathsf T}\mathbf A_n(a)\) has second off-diagonals equal to \(a\).
\end{itemize}
\end{proof}

\begin{corollary}[Toeplitz reduction]\label{cor:Toeplitz-reduction}
Let \(\Delta_0:=1\), and for \(m\ge 1\) define
\[
\Delta_m:=\det \mathbf\Theta_m(d,c,a),
\]
where
\[
\mathbf\Theta_1(d,c,a):=[d],
\qquad
\mathbf\Theta_2(d,c,a):=
\begin{bmatrix}
d & -c\\
-\overline c & d
\end{bmatrix},
\]
and for \(m\ge 3\),
\[
\mathbf\Theta_m(d,c,a)=
\begin{bmatrix}
d & -c & a \\
-\overline c & d & -c & \ddots \\
a & -\overline c & d & \ddots & \ddots \\
& \ddots & \ddots & \ddots & \ddots & \ddots \\
&& \ddots & \ddots & d & -c & a \\
&&& \ddots & -\overline c & d & -c \\
&&&& a & -\overline c & d
\end{bmatrix}.
\]
Then
\[
F_{n,\varepsilon}(x,y;a)
=
\Delta_n-(a^2+1)\Delta_{n-1}+a^2\Delta_{n-2}.
\]
\end{corollary}

\begin{proof}
By Proposition~\ref{prop:explicit-H},
\[
\mathbf H_{n,\varepsilon}(x,y;a)
=
\mathbf\Theta_n(d,c,a)-\mathbf E_{11}-a^2\mathbf E_{nn},
\]
up to the obvious coefficients matching the first and last diagonal entries. Since the determinant is affine in each of the modified diagonal entries, deleting the first and/or last row and column yields
\[
F_{n,\varepsilon}
=
\Delta_n-(a^2+1)\Delta_{n-1}+a^2\Delta_{n-2}.
\]
\end{proof}

\begin{remark}[Closed form]\label{rem:quartic}
Because \(\mathbf\Theta_m(d,c,a)\) is banded Toeplitz, the determinants \(\Delta_m\) satisfy a fourth-order linear recurrence in \(m\). Equivalently, \(\Delta_m\) is a linear combination of powers of the roots of the reciprocal quartic
\[
a\rho^4-c\rho^3+d\rho^2-\overline c\,\rho+a=0.
\]
Thus the boundary polynomial admits a closed form in \(n\), even though that form is not especially illuminating for the topology.
\end{remark}

\section{Vertical slices and the slice index}

We first record a general criterion that reduces simple connectedness to a one-dimensional statement along vertical lines.

\begin{theorem}[Vertical-slice criterion]\label{thm:vertical-slice}
Fix \(n\ge 2\), \(0<a<1\), and \(\varepsilon>0\). Assume that for every connected component \(\Omega\) of \(\sigma_\varepsilon(\mathbf A_n(a))\) and every \(x\in\R\), the set
\[
\{y\ge 0:\ x+iy\in \partial\Omega\}
\]
has at most one element. Then every connected component of \(\sigma_\varepsilon(\mathbf A_n(a))\) is simply connected.
\end{theorem}

\begin{proof}
Let \(\Omega\) be a connected component of \(\sigma_\varepsilon(\mathbf A_n(a))\). Since \(\|(z\mathbf I_n-\mathbf A_n(a))^{-1}\|\to 0\) as \(|z|\to\infty\), the set \(\Omega\) is bounded.

We claim that \(\Omega\) contains an eigenvalue of \(\mathbf A_n(a)\). If not, then \(\overline\Omega\cap \sigma(\mathbf A_n(a))=\varnothing\), so the resolvent norm
\[
\nu(z):=\|(z\mathbf I_n-\mathbf A_n(a))^{-1}\|
\]
is continuous on \(\overline\Omega\) and subharmonic in a neighborhood of \(\overline\Omega\). By Proposition~\ref{prop:boundary-polynomial},
\[
\nu(z)=\varepsilon^{-1}\quad (z\in\partial\Omega),
\]
whereas \(\nu(z)>\varepsilon^{-1}\) on \(\Omega\). This contradicts the maximum principle for subharmonic functions. Hence \(\Omega\cap \sigma(\mathbf A_n(a))\neq\varnothing\). By Lemma~\ref{lem:real-spectrum}, every eigenvalue is real, so \(\Omega\cap\R\neq\varnothing\).

Because \(\mathbf A_n(a)\) is real,
\[
\overline{\sigma_\varepsilon(\mathbf A_n(a))}=\sigma_\varepsilon(\mathbf A_n(a)),
\]
and therefore the complex conjugate
\[
\Omega^\sharp:=\{\overline z:\ z\in\Omega\}
\]
is also a connected component. Since \(\Omega\cap\R\neq\varnothing\), the components \(\Omega\) and \(\Omega^\sharp\) intersect, hence \(\Omega=\Omega^\sharp\). Thus every vertical fiber
\[
\Omega_x:=\{y\in\R:\ x+iy\in\Omega\}
\]
is symmetric with respect to \(y\mapsto -y\).

Fix \(x\in\R\). The set \(\Omega_x\) is an open bounded subset of \(\R\), symmetric under \(y\mapsto -y\). By hypothesis, its nonnegative boundary has at most one point. Therefore either \(\Omega_x=\varnothing\) or
\[
\Omega_x=(-h(x),h(x))
\]
for some \(h(x)>0\). Define
\[
I:=\{x\in\R:\ \Omega_x\neq\varnothing\}.
\]
Then \(I\) is open, and
\[
\Omega=\{x+iy:\ x\in I,\ |y|<h(x)\}.
\]
If \(I\) were disconnected, then \(\Omega\) would split accordingly, contradicting connectedness. Hence \(I\) is an interval. The deformation
\[
\Phi_t(x+iy):=x+i(1-t)y,
\qquad 0\le t\le 1,
\]
retracts \(\Omega\) onto \(I\). Thus \(\Omega\) is contractible, hence simply connected.
\end{proof}

A convenient scalar indicator for vertical slices is the negative index of the Hermitian boundary pencil.

\begin{proposition}[Slice index criterion]\label{prop:slice-index}
For fixed \(x\in\R\), define
\[
N_x(y):=\ind \mathbf H_{n,\varepsilon}(x,y;a),
\qquad y\ge 0.
\]
Then
\[
x+iy\in \sigma_\varepsilon(\mathbf A_n(a))
\iff
N_x(y)\ge 1.
\]
Equivalently,
\[
\{y\ge 0:\ x+iy\in \sigma_\varepsilon(\mathbf A_n(a))\}
=
\{y\ge 0:\ N_x(y)\ge 1\}.
\]
\end{proposition}

\begin{proof}
The eigenvalues of \(\mathbf H_{n,\varepsilon}(x,y;a)\) are
\[
s_j((x+iy)\mathbf I_n-\mathbf A_n(a))^2-\varepsilon^2,
\qquad j=1,\dots,n.
\]
Hence \(\mathbf H_{n,\varepsilon}(x,y;a)\) has a negative eigenvalue if and only if some singular value of \((x+iy)\mathbf I_n-\mathbf A_n(a)\) is strictly smaller than \(\varepsilon\), which is equivalent to \(x+iy\in \sigma_\varepsilon(\mathbf A_n(a))\).
\end{proof}

\begin{corollary}[Index monotonicity implies simple connectedness]\label{cor:index-monotone}
Assume that for every \(x\in\R\), the function \(y\mapsto N_x(y)\) is nonincreasing on \([0,\infty)\). Then every connected component of \(\sigma_\varepsilon(\mathbf A_n(a))\) is simply connected.
\end{corollary}

\begin{proof}
By Proposition~\ref{prop:slice-index},
\[
I_x^+:=\{y\ge 0:\ x+iy\in \sigma_\varepsilon(\mathbf A_n(a))\}
=
\{y\ge 0:\ N_x(y)\ge 1\}.
\]
Since \(N_x\) is integer-valued and nonincreasing, \(I_x^+\) is an initial segment of \([0,\infty)\). Because the pseudospectrum is open, either \(I_x^+=\varnothing\) or \(I_x^+=[0,h_x)\) for some \(h_x>0\). By conjugation symmetry, the full slice is either empty or
\[
(-h_x,h_x).
\]
Thus the hypothesis of Theorem~\ref{thm:vertical-slice} is satisfied.
\end{proof}

\begin{remark}
Corollary~\ref{cor:index-monotone} is the natural point at which discrete oscillation theory and matrix Prüfer methods enter; compare \cite{Hilscher2012,KratzHilscher2013,SchulzBaldes2007}. The final section below takes a different route: rather than proving global slice monotonicity directly, it rules out hole formation componentwise on the remaining compact parameter strip.
\end{remark}

\section{A componentwise topological proof sketch on the remaining strip}

In this section we recast the final step in a normalized chart adapted to the regular complex-root branch. The argument is explicitly a \emph{proof sketch}: it packages the local differential geometry and the semialgebraic continuation argument in a way that cleanly excludes hole creation. The purpose is to complement, rather than replace, the analytic reductions above.

For fixed \(n\ge 2\) and \(\varepsilon>0\), set
\[
P_{a,\varepsilon}:=\sigma_\varepsilon(\mathbf A_n(a)),
\qquad 0<a\le 1.
\]

\begin{proposition}[Standing input from the endpoint regimes]\label{prop:standing-input}
Fix \(n\ge 2\) and \(\varepsilon>0\). There exist numbers
\[
0<a_-<\widehat a_n<1
\]
such that the following hold.

\begin{enumerate}
\item For every \(0<a<a_-\), the set \(P_{a,\varepsilon}\) is a single star-shaped Jordan domain. In particular, every connected component of \(P_{a,\varepsilon}\) is simply connected.

\item For every \(\widehat a_n<a\le 1\), every connected component of \(P_{a,\varepsilon}\) is simply connected.

\item On the compact strip
\[
I:=[a_-,\widehat a_n],
\]
every boundary point of \(P_{a,\varepsilon}\) lies on the regular complex-root branch. In particular, on \(I\) there are no real-root boundary points and no other nongeneric boundary points outside the regular complex-root chart.
\end{enumerate}
\end{proposition}

\begin{remark}
Proposition~\ref{prop:standing-input} is the only nonlocal input in the final argument. It is consistent with the two endpoint models of Remark~\ref{rem:endpoints}: \(a=0\) is the Jordan-block limit, while \(a=1\) is the normal limit.
\end{remark}

\subsection*{The regular complex-root chart}

Fix \(a\in I\), and put
\[
\tau:=1-a>0.
\]
Introduce the upper-half-plane variables
\[
u:=x^2,\qquad v:=y^2,
\]
and define
\[
\beta(u,v):=\frac{u+v-\varepsilon^2+\tau^2}{a},
\qquad
t(u,v):=\frac{(1+a)^2u+\tau^2v}{4a^2}.
\]
Let
\[
\mathbf T_n:=\tridiag(1,0,1),
\qquad
\mathbf E_{nn}:=(\delta_{jn}\delta_{kn})_{j,k=1}^n,
\qquad
\mathbf e_1:=(1,0,\dots,0)^{\mathsf T}.
\]
On the regular complex-root branch one arrives at the \(n\times n\) Hermitian matrix
\[
\mathbf N_a(t,\beta)
:=
\mathbf T_n^2-2\sqrt t\,\mathbf T_n+\beta\mathbf I_n+\tau\mathbf E_{nn}
=
(\mathbf T_n-\sqrt t\,\mathbf I_n)^2+(\beta-t)\mathbf I_n+\tau\mathbf E_{nn}.
\]
On the branch under consideration one has \(\mathbf N_a(t,\beta)\succ 0\), so the scalar function
\[
m_n(t,\beta;a)
:=
\mathbf e_1^{\mathsf T}\mathbf N_a(t,\beta)^{-1}\mathbf e_1
\]
is well defined. The boundary equation on the upper-half-plane chart is then
\[
\Phi_a(u,v)
:=
m_n(t(u,v),\beta(u,v);a)-\frac{a}{\tau}
=
0.
\]
We also set
\[
\widetilde\Phi_a(x,v):=\Phi_a(x^2,v),
\qquad v\ge 0.
\]

\begin{remark}\label{rem:rational}
By the adjugate formula, the entries of \(\mathbf N_a(t,\beta)^{-1}\) are rational functions of \(\sqrt t\), \(\beta\), and \(a\). Hence \(m_n\), \(\Phi_a\), and \(\widetilde\Phi_a\) are real-analytic and semialgebraic on the regular complex-root chart.
\end{remark}

\begin{proposition}[Monotonicity input on the regular chart]\label{prop:monotone-m}
For every \(a\in I\) and every regular complex-root point,
\[
\partial_t m_n(t,\beta;a)>0,
\qquad
\partial_\beta m_n(t,\beta;a)<0.
\]
Moreover,
\[
\partial_\beta m_n(t,\beta;a)
=
-\mathbf e_1^{\mathsf T}\mathbf N_a(t,\beta)^{-2}\mathbf e_1<0.
\]
\end{proposition}

\begin{lemma}[Derivative identity]\label{lem:derivative-identity}
For every \(a\in I\),
\[
\partial_u\Phi_a
=
\frac1a\,\partial_\beta m_n
+
\frac{(1+a)^2}{4a^2}\,\partial_t m_n,
\]
\[
\partial_v\Phi_a
=
\frac1a\,\partial_\beta m_n
+
\frac{\tau^2}{4a^2}\,\partial_t m_n.
\]
Hence
\[
\partial_u\Phi_a-\partial_v\Phi_a
=
\frac1a\,\partial_t m_n
>0.
\]
\end{lemma}

\begin{proof}
This is a direct chain-rule computation using
\[
\partial_u\beta=\partial_v\beta=\frac1a,
\qquad
\partial_u t=\frac{(1+a)^2}{4a^2},
\qquad
\partial_v t=\frac{\tau^2}{4a^2}.
\]
Subtracting the two identities gives
\[
\partial_u\Phi_a-\partial_v\Phi_a
=
\frac{(1+a)^2-\tau^2}{4a^2}\,\partial_t m_n
=
\frac{4a}{4a^2}\,\partial_t m_n
=
\frac1a\,\partial_t m_n>0.
\]
\end{proof}

\begin{lemma}[No off-axis critical points]\label{lem:no-off-axis-critical}
Fix \(a\in I\). Let \(x\neq 0\) and \(v>0\). If
\[
\widetilde\Phi_a(x,v)=0,
\]
then
\[
\nabla \widetilde\Phi_a(x,v)\neq 0.
\]
\end{lemma}

\begin{proof}
We have
\[
\partial_x\widetilde\Phi_a(x,v)=2x\,\partial_u\Phi_a(x^2,v),
\qquad
\partial_v\widetilde\Phi_a(x,v)=\partial_v\Phi_a(x^2,v).
\]
If both vanished, then \(x\neq 0\) would force
\[
\partial_u\Phi_a(x^2,v)=0,
\qquad
\partial_v\Phi_a(x^2,v)=0,
\]
contradicting Lemma~\ref{lem:derivative-identity}.
\end{proof}

\begin{lemma}[Real-axis tangencies]\label{lem:real-axis-tangency}
Fix \(a\in I\), and let \(x_0\neq 0\). Suppose
\[
\widetilde\Phi_a(x_0,0)=0,
\qquad
\partial_x\widetilde\Phi_a(x_0,0)=0.
\]
Then
\[
\partial_v\widetilde\Phi_a(x_0,0)<0.
\]
Consequently, near \((x_0,0)\), the zero set of \(\widetilde\Phi_a\) is the graph
\[
v=\psi(x),
\qquad
\psi(x_0)=0,
\qquad
\psi'(x_0)=0.
\]
Equivalently, in \((x,y)\)-coordinates the boundary is locally of the form
\[
y^2=\psi(x).
\]
\end{lemma}

\begin{proof}
Since \(x_0\neq 0\),
\[
0=\partial_x\widetilde\Phi_a(x_0,0)=2x_0\,\partial_u\Phi_a(x_0^2,0),
\]
hence \(\partial_u\Phi_a(x_0^2,0)=0\). By Lemma~\ref{lem:derivative-identity},
\[
0<\partial_u\Phi_a(x_0^2,0)-\partial_v\Phi_a(x_0^2,0),
\]
so \(\partial_v\Phi_a(x_0^2,0)<0\), i.e.
\[
\partial_v\widetilde\Phi_a(x_0,0)<0.
\]
The implicit-function theorem now gives the stated graph \(v=\psi(x)\), and differentiating the identity \(\widetilde\Phi_a(x,\psi(x))=0\) at \(x=x_0\) yields \(\psi'(x_0)=0\).
\end{proof}

\begin{lemma}[Imaginary-axis tangencies]\label{lem:imag-axis-tangency}
Fix \(a\in I\), and let \(v_0>0\). Suppose
\[
\widetilde\Phi_a(0,v_0)=0,
\qquad
\partial_v\widetilde\Phi_a(0,v_0)=0.
\]
Then
\[
\partial_u\Phi_a(0,v_0)>0.
\]
Consequently, near \((0,v_0)\) in \((u,v)\)-coordinates the boundary is a graph
\[
u=\phi(v),
\qquad
\phi(v_0)=0,
\qquad
\phi'(v_0)=0.
\]
If \(y_0:=\sqrt{v_0}\), then near \((0,\pm y_0)\) in \((x,y)\)-coordinates the boundary is locally of the form
\[
x^2=\chi_\pm(y),
\qquad
\chi_\pm(\pm y_0)=0,
\qquad
\chi_\pm'(\pm y_0)=0.
\]
\end{lemma}

\begin{proof}
At \(x=0\), one automatically has \(\partial_x\widetilde\Phi_a(0,v)=0\). The hypothesis gives
\[
0=\partial_v\widetilde\Phi_a(0,v_0)=\partial_v\Phi_a(0,v_0).
\]
By Lemma~\ref{lem:derivative-identity},
\[
\partial_u\Phi_a(0,v_0)>\partial_v\Phi_a(0,v_0)=0.
\]
Hence \(\partial_u\Phi_a(0,v_0)>0\), and the implicit-function theorem yields the graph \(u=\phi(v)\) with \(\phi(v_0)=0\). Differentiating \(\Phi_a(\phi(v),v)=0\) at \(v=v_0\) gives \(\phi'(v_0)=0\). Passing from \(v=y^2\) back to \(y\) yields the two local models \(x^2=\chi_\pm(y)\) near \((0,\pm y_0)\).
\end{proof}

\begin{remark}[Local fold models]\label{rem:fold}
Lemmas~\ref{lem:real-axis-tangency} and \ref{lem:imag-axis-tangency} show that every axis singularity on the remaining strip is of fold type. In the \((x,y)\)-plane these may appear as an eye-shaped cap, a cusp, or two symmetric branches meeting on an axis. The only fact needed below is that these are \emph{minimum-type} axis contacts; they are not saddle events.
\end{remark}

\subsection*{Local topology of the fold models}

\begin{proposition}[Axis-tangency models are hole-free]\label{prop:axis-hole-free}
Let \(R\subset\R^2\) be a sufficiently small closed rectangle.

\begin{enumerate}
\item Suppose
\[
R\cap \partial\Omega
=
\{(x,y)\in R:\ y^2=\psi(x)\},
\]
where \(\psi\in C^1\) and \(\psi(x_0)=\psi'(x_0)=0\) at some interior point \(x_0\). Then every connected component of \(R\cap\Omega\) is simply connected, and every connected component of \(R\setminus\Omega\) meets \(\partial R\).

\item Suppose
\[
R\cap \partial\Omega
=
\{(x,y)\in R:\ x^2=\chi(y)\},
\]
where \(\chi\in C^1\) and \(\chi(y_0)=\chi'(y_0)=0\) at some interior point \(y_0\). Then the same conclusion holds.
\end{enumerate}

In particular, neither local model can create a new bounded complementary component inside \(R\).
\end{proposition}

\begin{proof}
We prove (1); the second case is identical after exchanging the roles of \(x\) and \(y\). After shrinking \(R\) if necessary, the set
\[
\{(x,y)\in R:\ y^2=\psi(x)\}
\]
contains no closed loop inside \(R\). Hence \(R\setminus \partial\Omega\) decomposes into regions of the form
\[
\{(x,y)\in R:\ y^2<\psi(x)\}
\quad\text{and}\quad
\{(x,y)\in R:\ y^2>\psi(x)\},
\]
each of which is a topological disc or an exterior strip touching the boundary of \(R\). Therefore every connected component of \(R\cap \Omega\) is simply connected, and every connected component of \(R\setminus\Omega\) touches \(\partial R\).
\end{proof}

\begin{proposition}[Fold singularities do not create holes]\label{prop:fold-no-hole}
Consider a semialgebraic one-parameter family of planar open sets. Suppose that at a critical parameter value \(a_*\) every local topology-changing point of the boundary is, after a local semialgebraic change of coordinates, of one of the two types covered by Proposition~\ref{prop:axis-hole-free}. Then crossing \(a=a_*\) may create or annihilate connected components, or merge or split already existing components, but it cannot create a new bounded complementary component of any connected component.
\end{proposition}

\begin{proof}
Proposition~\ref{prop:axis-hole-free} shows that each such local model is hole-free in a sufficiently small rectangle: locally, both the set and its complement are unions of topological discs, and no component of the complement is trapped inside the rectangle. Thus the only possible local surgeries are \(0\)-handle events (birth/death of a disc or merger/splitting along a disc attachment). In the plane, creation of a hole would require a \(1\)-handle (saddle-type) event. Since saddle models are excluded, hole birth is impossible.
\end{proof}

\subsection*{Topological preliminaries}

\begin{lemma}[Planar criterion]\label{lem:planar-criterion}
A connected open set \(\Omega\subset\R^2\) is simply connected if and only if \(\R^2\setminus\Omega\) has no bounded connected component.
\end{lemma}

\begin{proof}
This is standard planar topology: after one-point compactification, \(\Omega\) is simply connected if and only if its complement in \(S^2\) is connected. Since the component containing \(\infty\) is always present, this is equivalent to saying that \(\R^2\setminus\Omega\) has no bounded connected component.
\end{proof}

\begin{lemma}[Uniform boundedness and semialgebraicity]\label{lem:semialg}
Fix \(\varepsilon>0\) and let \(J\subset (0,1]\) be compact. Then
\[
\mathcal X_J
:=
\{(x,y,a)\in \R^2\times J:\ x+iy\in P_{a,\varepsilon}\}
\]
is bounded and semialgebraic.
\end{lemma}

\begin{proof}
Write
\[
\mathbf A_n(a)=\mathbf U+a\mathbf L,
\]
where \(\mathbf U\) is the upper shift and \(\mathbf L\) is the lower shift. Since \(\|\mathbf U\|_2=\|\mathbf L\|_2=1\) and \(0<a\le 1\),
\[
\|\mathbf A_n(a)\|_2\le 2.
\]
If \(z\in P_{a,\varepsilon}\), then
\[
s_{\min}(z\mathbf I_n-\mathbf A_n(a))<\varepsilon,
\]
but also
\[
s_{\min}(z\mathbf I_n-\mathbf A_n(a))
\ge |z|-\|\mathbf A_n(a)\|_2
\ge |z|-2.
\]
Hence \(|z|<2+\varepsilon\), uniformly for \(a\in J\). This proves boundedness.

For semialgebraicity, write \(z=x+iy\) and \(\boldsymbol\xi=\mathbf u+i\mathbf v\) with \(\mathbf u,\mathbf v\in\R^n\). Then
\[
(x,y,a)\in \mathcal X_J
\]
if and only if there exist \(\mathbf u,\mathbf v\in\R^n\) such that
\[
\|\mathbf u\|^2+\|\mathbf v\|^2=1
\]
and
\[
\|((x\mathbf I_n-\mathbf A_n(a))\mathbf u-y\mathbf v)\|^2
+
\|(y\mathbf u+(x\mathbf I_n-\mathbf A_n(a))\mathbf v)\|^2
<\varepsilon^2.
\]
These are polynomial equalities and inequalities in real variables, so \(\mathcal X_J\) is semialgebraic by the Tarski--Seidenberg theorem.
\end{proof}

\begin{lemma}[Local triviality away from critical events]\label{lem:local-trivial}
Fix \(a_*\in I\), and let \(p\in \partial P_{a_*,\varepsilon}\). Assume \(p\) lies on the regular complex-root branch and is not one of the axis-tangency points of Lemmas~\ref{lem:real-axis-tangency} and \ref{lem:imag-axis-tangency}. Then there exist a neighborhood \(U\subset\R^2\) of \(p\) and \(\delta>0\) such that the pairs
\[
\bigl(U,U\cap P_{a,\varepsilon}\bigr),
\qquad
|a-a_*|<\delta,
\]
are mutually homeomorphic.
\end{lemma}

\begin{proof}
If \(p\) is off the symmetry axes, Lemma~\ref{lem:no-off-axis-critical} and the implicit-function theorem show that the boundary is a smooth embedded arc depending smoothly on \(a\). If \(p=(x_0,0)\) with \(x_0\neq 0\) and \(\partial_x\widetilde\Phi_{a_*}(x_0,0)\neq 0\), then the implicit-function theorem writes the boundary locally as \(x=\chi(v,a)\). If \(p=(0,\pm y_0)\) with \(y_0>0\) and \(\partial_v\widetilde\Phi_{a_*}(0,y_0^2)\neq 0\), then locally one has \(v=\nu(x,a)\), equivalently an embedded arc in \((x,y)\)-coordinates. In each case, the local side decomposition is constant for \(a\) near \(a_*\).
\end{proof}

\subsection*{The final theorem}

\begin{theorem}[Componentwise simple connectedness; proof sketch]\label{thm:main}
For every \(n\ge 2\), every \(0<a<1\), and every \(\varepsilon>0\), every connected component of
\[
P_{a,\varepsilon}
=
\sigma_\varepsilon(\mathbf A_n(a))
\]
is simply connected.

Consequently, by Remark~\ref{rem:general-b}, every connected component of \(\sigma_\varepsilon(\tridiag(\alpha,0,\beta))\) is simply connected whenever \(0<\alpha<\beta\).
\end{theorem}

\begin{proof}[Proof sketch]
By Proposition~\ref{prop:standing-input}, the statement is already known in the low-\(a\) regime \(0<a<a_-\) and in the near-normal regime \(\widehat a_n<a\le 1\). Thus only the compact strip
\[
I=[a_-,\widehat a_n]
\]
remains.

Choose \(a_0\in(0,a_-)\) and consider
\[
\mathcal X
:=
\{(x,y,a)\in \R^2\times [a_0,\widehat a_n]:\ x+iy\in P_{a,\varepsilon}\}.
\]
By Lemma~\ref{lem:semialg}, \(\mathcal X\) is bounded and semialgebraic. Hence Hardt's semialgebraic triviality theorem (see, e.g., \cite{BCR}) implies that the projection
\[
\pi:\mathcal X\to [a_0,\widehat a_n],\qquad (x,y,a)\mapsto a,
\]
has only finitely many exceptional parameter values. On each component of the complement of that finite set, the fibers \(P_{a,\varepsilon}\) are semialgebraically homeomorphic.

Suppose, for contradiction, that the theorem fails. Since all fibers with \(a<a_-\) have simply connected components, there exists a first exceptional value
\[
a_*\in I
\]
such that for \(a>a_*\) sufficiently close to \(a_*\), some connected component of \(P_{a,\varepsilon}\) is not simply connected, whereas for \(a<a_*\) sufficiently close to \(a_*\), every connected component is simply connected.

By Lemma~\ref{lem:planar-criterion}, this means that when \(a\) crosses \(a_*\) from left to right, some connected component acquires a new bounded complementary component, that is, a new hole.

We now classify the only possible local topology-changing points on \(\partial P_{a_*,\varepsilon}\).

\begin{itemize}
\item By Proposition~\ref{prop:standing-input}(3), every boundary point on the remaining strip lies on the regular complex-root branch.

\item If such a point lies off the symmetry axes, then Lemma~\ref{lem:no-off-axis-critical} shows that it is regular; hence Lemma~\ref{lem:local-trivial} gives local topological constancy.

\item If the point lies on the real axis away from the origin, the only possible singularities are the fold-type tangencies of Lemma~\ref{lem:real-axis-tangency}, with local model \(y^2=\psi(x)\).

\item If the point lies on the imaginary axis away from the origin, the only possible singularities are the fold-type tangencies of Lemma~\ref{lem:imag-axis-tangency}, with local model \(x^2=\chi(y)\).
\end{itemize}

Thus \emph{every} local topology-changing point at \(a=a_*\) is a fold singularity of one of the two axis types covered by Proposition~\ref{prop:axis-hole-free}. By Proposition~\ref{prop:fold-no-hole}, such singularities may create or destroy components, or merge or split already existing components, but they cannot create a bounded complementary component.

This contradicts the definition of \(a_*\) as the first parameter where some connected component ceases to be simply connected. Therefore no such bad parameter exists on \(I\). Combining this with the two endpoint regimes from Proposition~\ref{prop:standing-input} proves the theorem.
\end{proof}

\begin{remark}
The proof is componentwise rather than global. Several disconnected pseudospectral components may occur. What the argument excludes is the birth of a \emph{hole inside a fixed connected component}.
\end{remark}

\begin{thebibliography}{99}

\bibitem{BCR}
J.~Bochnak, M.~Coste, and M.-F.~Roy,
\newblock \emph{Real Algebraic Geometry},
\newblock Ergebnisse der Mathematik und ihrer Grenzgebiete (3), Vol.~36,
Springer, Berlin, 1998.

\bibitem{Hilscher2012}
R.~{\v S}imon Hilscher,
\newblock Oscillation theorems for discrete symplectic systems with nonlinear dependence in spectral parameter,
\newblock \emph{Linear Algebra Appl.} \textbf{437} (2012), no.~12, 2922--2960.

\bibitem{KratzHilscher2013}
W.~Kratz and R.~{\v S}imon Hilscher,
\newblock A generalized index theorem for monotone matrix-valued functions with applications to discrete oscillation theory,
\newblock \emph{SIAM J. Matrix Anal. Appl.} \textbf{34} (2013), no.~1, 228--243.

\bibitem{SchulzBaldes2007}
H.~Schulz-Baldes,
\newblock Rotation numbers for Jacobi matrices with matrix entries,
\newblock \emph{Math. Phys. Electron. J.} \textbf{13} (2007), Paper No.~5, 40~pp.

\bibitem{TrefethenEmbree2005}
L.~N.~Trefethen and M.~Embree,
\newblock \emph{Spectra and Pseudospectra: The Behavior of Nonnormal Matrices and Operators},
\newblock Princeton University Press, Princeton, NJ, 2005.

\end{thebibliography}

\end{document}